\scanpage{563}
\mspage{585}

The kernel of $\varepsilon_{\mathcal{G}}$ is none other than $\mathcal{G}^{+}$. In \unsure{its capacity} as a subgroup of $\mathcal M_{\rho,\sigma}$ (resp. of $G_{\rho,\sigma}$) $\mathcal{G}$ is identified with the inverse image in this group of the subgroup $\{1,\tau'_{\Gamma}\}$ of $\Gamma_{\rho,\sigma}$ (resp. \inserted{of the subgroup} $\{1,\tau'_{\gamma}\}$ of $\gamma_{\rho,\sigma}$), where $\tau'_{\Gamma},\tau'_{\gamma}$ denote the images of $\tau'\in\mathcal{G}$ in $\Gamma_{\rho,\sigma}$ and $\gamma_{\rho,\sigma}$ respectively.

g) The diagram (62) has as its main part
\[
\text{(67)}\qquad
\diagram{\begin{tikzcd}[row sep=0.9em,column sep=1.25em]
{} & {}\arrow[rrrr]\arrow[ddr] & {} & {} & {} & {}\arrow[dd] & {} & {} & {} & {}\\
{} & {} & {} & {}\arrow[drr] & {} & {} & {} & {} & {} & {}\\
{}\arrow[uur]\arrow[ddr] & {} & {}\arrow[ur]\arrow[dr] & {} & {} & {}\arrow[urrr,dashed,shorten <=3em]\arrow[rrr,shorten <=3em]\arrow[drrr,dashed,shorten <=3em] & \Biggl[ & {} & {} & \Biggr]\\
{} & {} & {} & {}\arrow[urr] & {} & {} & {} & {} & {} & {}\\
\deleted{\unsure{[\mathcal{G}^{+}]}} & {}\arrow[uur]\arrow[rrrr] & {} & {} & {} & {}\arrow[uu] & {} & {} & {} & {}\\
\unsure{\Bigl[}{}\arrow[ur,dashed,hook]\arrow[urrrrr,dashed,hook] & {} & {} & {} & {} & \Bigr] & {} & {} & {} & {}
\end{tikzcd}},
\]
all the groups marked on this \inserted{last} diagram, except for $\widehat{\mathbb Z}^{*},\{\pm1\},\{\pm1\}$ (put in brackets $[\;]$) map to $\gamma_{\rho,\sigma}$, which is the terminal term of this diagram, linking \unsure{9} groups (not \unsure{counting} $\mathcal{G}\ \deleted{\unsure{\mathcal{G}^{+}}}$ inserted \deleted{\unsure{by}} as \deleted{\unsure{terms}}\inserted{\unsure{outside}} \deleted{\unsure{initial} of}\inserted{of two} dotted arrows) including $\gamma_{\rho,\sigma}$, obtained from four of them $Z_{\rho,\sigma},N_{\rho}^{*},N_{\sigma}^{*},G_{\rho,\sigma}$ (which map to $\gamma_{\rho,\sigma}$ by epimorphisms) by taking fiber products over $\gamma_{\rho,\sigma}$ of certain of them --- \deleted{\unsure{the}} the initial term being $S_{0}\mathcal M_{\rho,\sigma}$, which is the fiber product of the four preceding factors. Taking the kernels of the homomorphisms from these \deleted{four} nine groups to $\gamma_{\rho,\sigma}$, we find a diagram of the same combinatorial type, linking

\mspage{562}

{\small
\begin{center}
562\\
\texttt{https://grothendieck.umontpellier.fr}\\
Université de Montpellier, UMR CNRS 5149, Case courrier 051, Place Eugène Bataillon\\
34095 Montpellier cedex 5 - France
\end{center}
}

\scanpage{564}

nine groups, \deleted{including $\{1\}=\operatorname{Ker}(\gamma_{\rho,\sigma}\longrightarrow\gamma_{\rho,\sigma})$,} formed from $L_{0}=SZ_{\rho,\sigma}$, $SN_{\rho}^{*},SN_{\sigma}^{*}$ and $\mathcal{G}^{+}=SG_{\rho,\sigma}$ and from certain products of these groups (including the product over the empty index set, namely $\{1\}=\operatorname{Ker}(\gamma_{\rho,\sigma}\xrightarrow{\mathrm{id}}\gamma_{\rho,\sigma})$), with as \deleted{\unsure{arrowisms|arrows}} canonical projections. We have inserted in (62) only \deleted{five} of these product groups, with the canonical injections into the fiber products from which they arise, \deleted{in} dispensing with including in it \unsure{also|also} $SN_{\rho}^{*},SN_{\sigma}^{*},SG_{\rho,\sigma}=\mathcal{G}^{+}$, and the canonical inclusions $SN_{\rho}^{*}\hookrightarrow N_{\rho}^{*}$, $SN_{\sigma}^{*}\hookrightarrow N_{\sigma}^{*}$, $\mathcal{G}^{+}=SG_{\rho,\sigma}\hookrightarrow G_{\rho,\sigma}$, and the \deleted{\unsure{p}} \inserted{complete} diagram of the projections of the partial products of these groups and of $L_{0}$, \deleted{including} \inserted{which therefore \unsure{lacks}} the projections onto these \inserted{three} groups, and the homomorphisms from these groups and from $L_{0}$ to the unit group.

\mspage{563}

{\small
\begin{center}
563\\
\texttt{https://grothendieck.umontpellier.fr}\\
Université de Montpellier, UMR CNRS 5149, Case courrier 051, Place Eugène Bataillon\\
34095 Montpellier cedex 5 - France
\end{center}
}

\scanpage{565}
\mspage{586}

6°) \uline{The set} $\Sigma_{\rho,\sigma},S\unsure{\mathcal{G}|g}_{\rho,\sigma}$.

We put
\[
\text{(68)}\qquad
\begin{gathered}
\Sigma_{\rho,\sigma}\subset\mathcal{G}^{+}\times\mathcal{G}^{+}\times\widehat{\mathbb Z}^{*}\times\{\pm1\}\times\{\pm1\}\\
\begin{array}{c}
\Vert\\[-0.4em]
\displaystyle
\left\{(\xi,\eta,\mu,\unsure{\varepsilon',\varepsilon})\;\middle|\;
\begin{array}{l}
\xi\rho\ \inserted{\mathcal{G}^{+}\text{-}}\text{conjugate to }\rho^{\varepsilon'}\\
\eta\sigma\ \mathcal{G}^{+}\text{-conjugate to }\sigma^{\varepsilon}\\
\xi=\eta\underbrace{\varepsilon_{1}^{\mu-1}}_{\rho_{1}^{\nu}\ \text{where }\nu=\frac{\mu-1}{2}}\quad\deleted{\unsure{\text{(with}}}
\end{array}\right\}
\end{array}
\end{gathered}
\]

\uline{N.B.} If $\rho$ not conjugate to $\rho^{-1}$ in $\mathcal{G}^{+}$, $\varepsilon'$ known when we know $\xi,\eta$ --- likewise if $\sigma$ not conjugate to $\sigma^{-1}$ (hence $\sigma^{2}\ne1$), $\varepsilon'$ is determined by $\xi,\eta$. \unsure{If this homomorphism} $\widehat{\mathbb Z}\longrightarrow L_{0}$ $(n\longmapsto\varepsilon_{0}^{n})$ is injective, then $\mu$ is determined by $\xi,\eta$, and the relation $\xi=\eta\varepsilon_{1}^{\mu-1}$ can also be written (without bringing in $\mu$)
\[
\text{(69)}\qquad \eta^{-1}\xi\in L_{1}.
\]

Let on the other hand
\[
\text{(70)}\qquad
\begin{gathered}
S\mathcal{G}_{\rho,\sigma}\subset\mathcal{G}\times\mathcal{G}\times\widehat{\mathbb Z}^{*}\\
\begin{array}{c}
\Vert\\[-0.4em]
\displaystyle\left\{(\alpha,\beta,\mu)\;\middle|\;
\underbrace{[\alpha,\rho]}_{\xi}=[\beta,\sigma]\varepsilon_{1}^{\mu-1}\right\}
\end{array}
\end{gathered}
\]

We have a map
\[
\text{(71)}\qquad
\left\{\diagram{\begin{tikzcd}[row sep=1.5em,column sep=3em]
S\mathcal{G}_{\rho,\sigma}\arrow[r] & \Sigma_{\rho,\sigma}\\
(\alpha,\beta,\mu)\arrow[r,mapsto] & (\xi,\eta,\mu,\varepsilon',\varepsilon)
\end{tikzcd}}\right.
\quad\text{with}\quad
\left\{\begin{array}{l}
\xi=[\alpha,\rho]\\
\eta=[\beta,\sigma]\\
\mu\ \text{the same}\\
\varepsilon'=\unsure{\varepsilon}_{\mathcal{G}}(\alpha)\\
\varepsilon=\varepsilon_{\mathcal{G}}(\beta)
\end{array}\right.
\]

\margin{\uline{N.B.} $\xi=\alpha\rho(\alpha)^{-1}\in\mathcal{G}^{+}$, since we have $\varepsilon_{\mathcal{G}}(\xi)=\unsure{\dot{\alpha}|\bar{\alpha}}\rho(\unsure{\dot{\alpha}|\bar{\alpha}})=1$ since $\rho$ operates trivially on $\mathcal{G}/\mathcal{G}^{+}\simeq\{\pm1\}$, for the same reason $\eta=\beta\sigma(\beta)^{\unsure{-1}}\in\mathcal{G}^{+}$.}

which is \inserted{obviously} \uline{surjective}.

We define on the other hand

\mspage{564}

{\small
\begin{center}
564\\
\texttt{https://grothendieck.umontpellier.fr}\\
Université de Montpellier, UMR CNRS 5149, Case courrier 051, Place Eugène Bataillon\\
34095 Montpellier cedex 5 - France
\end{center}
}

\scanpage{566}

\[
\text{(72)}\qquad
\begin{gathered}
\Sigma_{\rho,\sigma}^{\mathrm{eff}}\subset\Sigma_{\rho,\sigma}\\
\begin{array}{c}
\Vert\\[-0.4em]
\displaystyle\left\{(\xi,\unsure{\eta},\mu,\varepsilon',\varepsilon)\in\Sigma_{\rho,\sigma}\;\middle|\;
\begin{array}{l}
\exists u\in\operatorname{Aut}(\mathcal{G}^{+})\ \text{with}\\
u(\rho)=\xi\rho,\quad u(\unsure{\sigma})=\eta\sigma
\end{array}\right\}
\end{array}
\end{gathered}
\]
\margin{i.e. $[u,\rho]=\xi$, $[u,\sigma]=\eta$}
(this $u$ will then be unique, \inserted{and \unsure{conversely}, \unsure{knowledge} of $u$ determines $\xi,\eta$.}) There is a $1$-$1$ correspondence between the elements of $\Sigma_{\rho,\sigma}^{\mathrm{eff}}$, and the systems $(u,\mu,\varepsilon',\varepsilon)\in\operatorname{Aut}(\mathcal{G}^{+})\times\widehat{\mathbb Z}^{*}\times\{\pm1\}\times\{\pm1\}$, such that we have \deleted{$u(\rho)\ \mathcal{G}^{+}\text{-conjugate}\unsure{\text{.}}$}
\[
\left\{\begin{array}{lll}
u(\rho)\text{ is}&\mathcal{G}^{+}\text{-conjugate to}&\rho^{\varepsilon'}\\
u(\sigma)&\rule{9em}{0.35pt}&\sigma^{\varepsilon}\\
u(\varepsilon_{0})=\varepsilon_{0}^{\mu}&\multicolumn{2}{l}{(\text{which expresses }\xi=\eta\rho_{1}^{\nu})}
\end{array}\right.
\]
--- now this set is none other than the subgroup $\unsure{Z_{\rho,\sigma}|\mathbb Z_{\rho,\sigma}}$ of $\operatorname{Aut}(\mathcal{G}^{+})\times\widehat{\mathbb Z}^{*}\times\{\pm1\}\times\{\pm1\}$. \inserted{(cf. (6) page\,577')}

We thus find a bijection
\[
\text{(73)}\qquad
\diagram{\begin{tikzcd}[row sep=2em,column sep=3em]
Z_{\rho,\sigma}\arrow[r,"\sim"] & \Sigma_{\rho,\sigma}^{\mathrm{eff}}\\
\begin{array}{c}\deleted{\unsure{(u,\mu,\varepsilon,\varepsilon')}}\\u\end{array}\arrow[r,mapsto] & (\xi,\eta,\mu,\varepsilon',\varepsilon)
\end{tikzcd}}
\quad\left\{\begin{array}{l}
\xi=[u,\rho]\\
\eta=[u,\sigma]\\
\mu=\chi(u)\\
\varepsilon'=\varepsilon'_{Z}(u)\\
\varepsilon=\varepsilon_{Z}(u)
\end{array}\right.\ .
\]

We denote by $S\mathcal{G}_{\rho,\sigma}^{\mathrm{eff}}$ the inverse image of $\Sigma_{\rho,\sigma}^{\mathrm{eff}}$ in $S\mathcal{G}_{\rho,\sigma}$ under (71), so we will have
\[
\text{(74)}\qquad
\begin{gathered}
S\mathcal{G}_{\rho,\sigma}^{\mathrm{eff}}\subset S\mathcal{G}_{\rho,\sigma}\\
\begin{array}{c}
\Vert\\[-0.4em]
\displaystyle\left\{\unsure{\{\alpha,\beta,\mu\}}\in\mathcal{G}\times\mathcal{G}\times\widehat{\mathbb Z}^{*}\;\middle|\;
\begin{array}{l}
[\alpha,\rho]=[\beta,\sigma]\varepsilon_{1}^{\mu-1}\\
\exists u\in\operatorname{Aut}(\mathcal{G}^{+})\ \text{such that we have}\\
\left\{\begin{array}{ll}
u(\rho)=\alpha(\rho)&(=\xi\rho,\text{ where }\xi=[\alpha,\rho])\\
u(\sigma)=\beta(\sigma)&(=\eta\sigma,\text{ where }\eta=[\beta,\sigma])
\end{array}\right.
\end{array}\right\}
\end{array}
\end{gathered}
\]
so that we have a \uline{surjective} \unsure{homomorphism} induced by (71)
\[
\text{(75)}\qquad
\left|\diagram{\begin{tikzcd}[column sep=3em]
S\mathcal{G}_{\rho,\sigma}^{\mathrm{eff}}\arrow[r]&\Sigma_{\rho,\sigma}^{\mathrm{eff}}
\end{tikzcd}}\right.\ .
\]

\mspage{565}

{\small
\begin{center}
565\\
\texttt{https://grothendieck.umontpellier.fr}\\
Université de Montpellier, UMR CNRS 5149, Case courrier 051, Place Eugène Bataillon\\
34095 Montpellier cedex 5 - France
\end{center}
}

\scanpage{567}
\mspage{587}

We have a canonical map
\[
\text{(76)}\qquad
\left\{\diagram{\begin{tikzcd}[row sep=2em,column sep=2em]
S_{0}\Gamma_{\rho,\sigma}=\unsure{Z|\mathbb Z}_{\rho,\sigma}\times_{\gamma_{\rho,\sigma}}N_{\rho}^{*}\times_{\gamma_{\rho,\sigma}}N_{\sigma}^{*}\arrow[r]&S\mathcal{G}_{\rho,\sigma}^{\mathrm{eff}}\\
x=\unsure{\{u,a,b\}|(u,a,b)}\arrow[r,mapsto] & (\alpha,\beta,\mu)
\end{tikzcd}}\right.
\quad\text{with}\quad
\left\{\begin{array}{l}
\alpha=u\overbrace{a^{-1}\tau'^{\unsure{\nu'(x)}}}^{\in Z_{\rho}}\\
\beta=u\underbrace{b^{-1}\tau'^{\nu(x)}}_{\in Z_{\sigma}}\\
\mu=\chi_{S\Gamma}(x)
\end{array}\right.
\]

\margin{to be more correct, one should write
\[\left\{\begin{array}{l}
\alpha=u\,i_{\rho}(a)^{-1}\tau'^{\nu'(x)}\\
\beta=u\,i_{\sigma}(b)^{-1}\tau'^{\nu(x)}
\end{array}\right.\]}

where we have defined $\nu'(x),\unsure{\nu(x)|\nu'(x)}\in\mathbb Z/2\mathbb Z$ by
\[
\deleted{\unsure{(}}\qquad
(-1)^{\nu'(x)}=\varepsilon'(x),\qquad(-1)^{\nu(x)}=\varepsilon(x).
\]

By definition of $S_{0}\Gamma_{\rho,\sigma}$ as a fiber product, we will indeed have $ua^{-1},ub^{-1}\in\mathcal{G}^{+}$, hence $\alpha,\beta\in\mathcal{G}$, we will moreover have
\[
\text{(77)}\qquad
\left\{\begin{array}{l}
\displaystyle\alpha(\rho)=u\left(\underbrace{a^{-1}\left(\underbrace{\tau'^{\nu'(x)}(\rho)}_{\rho^{\varepsilon'(x)}}\right)}_{(\rho^{\varepsilon'(x)})^{\varepsilon'(x)}=\rho}\right)=u(\rho),\quad\text{and likewise}\\[2em]
\beta(\sigma)=u(\sigma),
\end{array}\right.
\qquad\text{i.e.}\quad
\begin{array}{l}
[\alpha,\rho]=[u,\rho]\\
[\beta,\sigma]=[u,\sigma]
\end{array}
\]

\margin{\uline{N.B.} We have, for \unsure{every} $a\in N_{\rho}^{*}$, if $\varepsilon_{\rho}(a)=(-1)^{\nu_{\rho}(a)}$, $a\tau'^{\nu_{\rho}(a)}\in Z_{\rho}$, hence $[ua\tau'^{\nu_{\rho}(a)},\rho]=[u,\rho]$. Likewise for $\sigma\ldots$}

\deleted{so that} moreover the condition
$\underbrace{[\alpha,\rho]}_{[u,\rho]}=\underbrace{[\beta,\sigma]}_{[u,\sigma]}\varepsilon_{1}^{\mu-1}$ \deleted{\unsure{here}}
merely expresses the condition that $u(\varepsilon_{0})=\varepsilon_{0}^{\mu}$ which enters into the very definition of $Z_{\rho,\sigma}$.

\deleted{Moreover} \inserted{I say that} the map in question \inserted{(76)} is injective, \deleted{for} It suffices to see that when $(\alpha,\beta,\mu)$ are known, we know $u$ --- \inserted{indeed} we \deleted{have \unsure{then}} \inserted{in terms of $(\alpha,\beta,\mu)$}
\[
\text{(78)}\qquad
\left\{\begin{array}{l}
\nu'(x)=\unsure{\nu|\varepsilon}_{\mathcal{G}}(\alpha),\quad
\nu(x)=\unsure{\nu|\varepsilon}_{\mathcal{G}}(\beta),\quad
\deleted{\unsure{\mu=}}\ \chi_{S\Gamma}(x)=\unsure{\chi}_{Z}(u)=\mu\\[0.6em]
\deleted{\text{whence we obtain}}\quad\text{whence}\quad
a=\tau'^{\unsure{\nu|\varepsilon}_{\mathcal{G}}(\alpha)}\alpha^{-1}u,\quad
b=\tau'^{\unsure{\nu|\varepsilon}_{\mathcal{G}}(\beta)}\beta^{-1}u
\end{array}\right.
\]
\deleted{\unsure{we}} so if $u$ is known, we also know $a$ and $b$ hence $x$. On the other hand to determine \inserted{in terms of $\alpha,\beta,\mu$,} $u\in Z_{\rho,\sigma}\subset\operatorname{Aut}(\mathcal{G}^{+})\times\widehat{\mathbb Z}^{*}\times\{\pm1\}\times\{\pm1\}$ it suffices, in addition to (78), to determine \deleted{\unsure{it}} \deleted{\unsure{as}} the element of $\operatorname{Aut}(\mathcal{G}^{+})$ which it defines, which is

\mspage{566}

{\small
\begin{center}
566\\
\texttt{https://grothendieck.umontpellier.fr}\\
Université de Montpellier, UMR CNRS 5149, Case courrier 051, Place Eugène Bataillon\\
34095 Montpellier cedex 5 - France
\end{center}
}